% =====================================================================================================
\section{The assumptions}\label{sec:assumptions}

Each assumption has a number and a fixed key; the key in parentheses is the label used to cite it. Table~\ref{tab:usage}
lists the statements that use each one.

\subsection{The data model}\label{sec:a-all}

\begin{assumption}{GL}{Gaussian linear model}
The data follow~\eqref{eq:model}, with a fixed noise variance $\sigs^2>0$ and noise independent of the design.
\end{assumption}

\begin{assumption}{DIM}{Dimension}
$2\le p\le n^A$ for a fixed $A\ge1$.
\end{assumption}

\begin{assumption}{SP}{Sparsity}
The true model size $\sz\ge1$ is fixed.
\end{assumption}

\subsection{The predictors and the signal}\label{sec:a-post}

\begin{assumption}{GR}{Gaussian rows}
The rows of $\Xr$ are i.i.d.\ $\Normal_p(m_n,\Sig)$, where $\Sig$ is a correlation matrix.
\end{assumption}

\begin{assumption}{SGR}{Sub-Gaussian rows on sparse directions}
The raw rows are i.i.d.\ copies of a random vector $x\in\R^{\praw}$ with mean $m_n$ and covariance matrix $\Sig$, a
correlation matrix, and, for a fixed $L>0$, $\E\exp\{\lambda u^\top(x-m_n)\}\le\exp\{L^2\lambda^2u^\top\Sig u/2\}$ for every
$\lambda\in\R$ and every $u\in\R^{\praw}$ with at most $5(A+1)(\sz+1)\log n$ nonzero entries. Gaussian rows \as{GR} are
the case $L=1$.
\end{assumption}

\begin{assumption}{FSb}{Fixed signals, raw scale}
$b_{\mathrm{lo}}\le\abs{\beta^*_j}\le b_{\mathrm{hi}}$ for every $j\in\Sz$, with fixed $0<b_{\mathrm{lo}}\le b_{\mathrm{hi}}<\infty$.
\end{assumption}

\begin{assumption}{FSt}{Fixed signals, standardised scale}
$\yr=\alpha\one+X\tstar+\varepsilon$ with $\varepsilon\sim\Normal_n(0,\sigs^2I_n)$, a fixed $\sigs^2>0$, noise independent
of $\Xr$, a deterministic $\tstar$, $\sz\ge1$ fixed, and $\theta_{\mathrm{lo}}\le\abs{\theta^*_j}\le\theta_{\mathrm{hi}}$ for
$j\in\Sz$, with fixed $0<\theta_{\mathrm{lo}}\le\theta_{\mathrm{hi}}$.
\end{assumption}

\begin{assumption}{EFS}{Population eigenvalue floor}
$\lmin(\Sigma_{n,S})\ge\varphi^2$ for every $S$ with $\abs{S}\le5(A+1)(\sz+1)\log n$, for a fixed $\varphi\in(0,1]$ (for
instance, $\lmin(\Sig)\ge\varphi^2$).
\end{assumption}

\subsection{The distribution-level setting}\label{sec:a-dist}

In the setting of Definition~\ref{def:dist-setting}, fix a \emph{level} $M\ge1$; the \emph{window} is
$q_n:=\min\{\praw,\lceil C_qM\log n\rceil\}$ with $C_q:=2(s(A+2)+4)/\log2$.

\begin{assumption}{dist}{Distribution-level setting; A6 Assumption~3}
The integer $s\ge0$ and the constants $A\ge1$, $L\ge1$, $c_\psi>0$ and $\bar b>0$ are fixed.
\begin{asmitems}
  \item\label{asm:dist-model} \emph{Gaussian linear model for i.i.d.\ pairs.} $y=\alpha^*+x^\top\bstar+\varepsilon$ with
    $\varepsilon\sim\Normal(0,\sigs^2)$, $\sigs^2>0$, independent of $x$; $\alpha^*$, $\bstar$, $\sigs$ and the law of $x$ may
    depend on $n$.
  \item\label{asm:dist-sparse} \emph{Sparsity and dimension.} At most $s$ coordinates of $\bstar$ are nonzero, and
    $2\le\praw\le n^A$.
  \item\label{asm:dist-subg} \emph{Sub-Gaussian predictors on sparse directions.} Every predictor has positive variance,
    and $\E\exp(\lambda u^\top\zeta)\le\exp\big(L^2\lambda^2\Var(u^\top\zeta)/2\big)$ for every $\lambda\in\R$ and every
    $u\in\R^{\praw}$ with at most $q_n$ nonzero entries.
  \item\label{asm:dist-nec} \emph{No near-exact collinearity.} $\lmin(\Corr(x_T))\ge c_\psi/n$ for every set $T$ of at
    most $q_n$ predictors.
  \item\label{asm:dist-bs} \emph{Bounded signal.} $\abs{\beta^*_j}\sd(x_j)\le\bar b\,\sigs$ for every $j$.
\end{asmitems}
\end{assumption}

\subsection{The reported-model assumptions}\label{sec:a-rm}

A \emph{leakage set} is a deterministic $Q$ with $\Sz\subseteq Q\subseteq\set{p}$; its \emph{lattice} is
$\Lat(Q):=\{T:\Sz\subseteq T\subseteq Q\}$. The notation $\Cov(\cdot,\cdot\mid W_U)$ and $\bar\beta_{U,l}$ is that of
Definition~\ref{def:designs}.

\begin{assumption}{V}{Visibility: no true predictor masked by $\le15$ others}
For every $j\in\Sz$ and every $U\not\ni j$ with $\abs{U}\le15$, $\abs{\Cov(W_j,W^\top\tstar\mid W_U)}\ge a$, for a fixed
$a>0$.
\end{assumption}

\begin{assumption}{reported}{Polynomial leakage set with local regressions; A6 Assumption~4}
Fix $b\in(0,\tfrac12)$ and $b_H>0$ with $15b_H+b<\tfrac12$.
\begin{asmitems}
  \item\label{asm:rr-leak} \emph{Leakage set.} There is a deterministic $Q\supseteq\Sz$ with $\abs{Q}\le n^b$ and, for
    large $n$, $\abs{Q}\le K$ and ($Q=\set{p}$ or $K\ge2\abs{Q}+32\sz+1$), such that: (i) \emph{no leak outside $Q$:}
    there is a deterministic $\varpi_n$ with $\varpi_n\sqrt{n/\log n}\to0$ such that
    $\abs{\Cov(W_l,W_j\mid W_U)}\le\varpi_n$ for all $l\notin Q$, $j\in\Sz\setminus U$ and $\abs{U}\le15$; and (ii)
    \emph{locality:} there is a deterministic $\delta_n\to0$, and every $l\notin Q$ has a deterministic
    $\Hl_l\supseteq\{l\}\cup\Sz$ with $\abs{\Hl_l}\le n^{b_H}$, such that for all $l\notin Q$, all $U\not\ni l$ with
    $\abs{U}\le15$ and $U_0:=U\cap\Hl_l$: $\abs{\Cov(W_l,W_v\mid W_{U_0})}\le\delta_n$ and
    $\abs{\Cov(W_j,W_v\mid W_{U_0})}\le\delta_n$ for all $v\in U\setminus\Hl_l$ and $j\in\Sz\setminus U_0$.
  \item\label{asm:rr-floor} \emph{Floor on the leakage set.} $\lmin(\Sigma_{n,Q\cup\{l\}})\ge\varphi^2$ for every
    $l\in\set{p}$, with $\varphi$ as in \as{EFS}.
  \item\label{asm:rr-local} \emph{Local regressions within the leakage set.} There are deterministic sequences
    $d_n=o(\log n)$ and $L_n$, and for every $l\in\set{p}$ a deterministic $N_l\subseteq Q\setminus\{l\}$ with
    $\abs{N_l}\le d_n$, such that $\norm{\bar\beta_{T,l}-\bar\beta_{T\cap N_l,l}}_1\le L_n$ for every $T\in\Lat(Q)$ with
    $l\notin T$, where $L_n\to0$ and $L_n^2\log(\abs{Q}+1)=o(\log n)$.
\end{asmitems}
\end{assumption}

\subsection{The certified-screening assumption}\label{sec:a-scr}

\begin{assumption}{SCR}{Floor on sets of the budget size, and signal bounds; A6 \S5.5}
The design is fixed. For a model-size budget $k_n\in\{0,\dots,K-1\}$: $\lmin(\Rhat_S)\ge\kappa$, for a fixed $\kappa>0$,
for every $S$ with $\abs{S}\le k_n+\sz$; $\tmin\ge\tlow>0$ and $\norm{\tstar}_1\le M_1$; $\kappa$, $\tlow$, $M_1$ and
$\sigs^2$ are fixed; and $(k_n+1)\log p+\log n\le c_{\mathrm{SB}}\,n$, where
$c_{\mathrm{SB}}:=\kappa^2\tlow^2/\{200(M_1^2+\sigs^2)\}$.
\end{assumption}

\subsection{The witness designs}\label{sec:a-wit}

\begin{assumption}{WIT}{Witness designs; A6 Theorem~6(b)}
$\sz$ is fixed; $\norm{\mu}^2\le nM_\mu$ and $\norm{\tstar}^2\le nM_\theta$ for fixed $M_\mu$, $M_\theta$; the null
predictors are either orthogonal (to each other and to the true ones) or have i.i.d.\ $\Normal(0,1)$ raw entries
independent of the true predictors and of the noise; and $p\le n^A$.
\end{assumption}

\subsection{The assumptions of the supplementary results}\label{sec:a-supp}

\begin{assumption}{BC}{Bounded coefficients}
$\max_{j}\abs{\theta^*_j}\le\bar\theta$ for a fixed $\bar\theta<\infty$.
\end{assumption}

\begin{assumption}{EF}{Eigenvalue floor on sets of bounded size}
$\lmin(\Rhat_S)\ge\varphi_0^2$ for every model $S$ with $\abs{S}\le2(A+1)\sz+2$, for a fixed $\varphi_0\in(0,1]$.
\end{assumption}

\begin{assumption}{NC}{No near-exact collinearity among sets of logarithmic size}
$\lmin(G_S)\ge g_0$ for every model $S$ with $\abs{S}\le3(A+1)(\sz+1)\log n$, for a fixed $g_0>0$.
\end{assumption}

\begin{assumption}{BM}{Beta-min}
Either (a) $\tmin\ge C_\beta\,\sigs\rn$ with a fixed $C_\beta>0$ such that $C_\beta^2\varphi_0^2>c_0(\sz,A)$, the constant
of Definition~\ref{def:constants}; or (b) $\tmin\sqrt{n/\log n}\to\infty$.
\end{assumption}

\begin{assumption}{margin}{Floor on the true model and an identifiability margin; A6 Assumption~A.4}
\begin{asmitems}
  \item $\lmin(\Rhat_{\Sz})\ge\varphi_0^2$ for a fixed $\varphi_0\in(0,1]$.
  \item Every model $S\in\Om$ that misses a true predictor and has at most $j_0:=\lceil(2A+1)\sz\rceil+1$ predictors
    outside $\Sz$ leaves unexplained signal $\norm{(I-P_S)\mu}^2\ge c_{\mathrm{id}}\,\sigs^2\log n$, for a fixed
    $c_{\mathrm{id}}>c_1(\sz,A)$, the constant of Definition~\ref{def:constants}.
\end{asmitems}
\end{assumption}

\begin{assumption}{omitted}{Candidate-only fitting; A6 Assumption~A.6}
The method is run on the response $\yr$ and a set $\mathcal C$ of $p$ candidate predictors, with rows $x_{\mathcal C,i}$
collected in $x_{\mathcal C}$; $\Sigma$ and $R$ are the covariance and the correlation matrix of a candidate row.
\begin{asmitems}
  \item \emph{Model.} The candidate rows are i.i.d.\ $\Normal_p(m_{\mathcal C},\Sigma)$, and, given them, the responses are
    independent with $\yr_i\mid x_{\mathcal C}\sim\Normal(\alpha'+x_{\mathcal C,i}^\top\beta^{\mathcal C},\sigma^2_{\mathcal C})$,
    $\sigma^2_{\mathcal C}>0$. We write $\Sz^{\mathcal C}:=\supp\beta^{\mathcal C}$ and $s_{\mathcal C}:=\abs{\Sz^{\mathcal C}}$.
  \item \emph{Floor.} $\lmin(R_T)\ge c_{\mathcal C}$ for every $T\subseteq\mathcal C$ with
    $\abs{T}\le3(A+1)(s_{\mathcal C}+1)\log n$, for a fixed $c_{\mathcal C}>0$.
  \item \emph{Signal relative to the residual variance.} $\beta^{\mathcal C\top}\Sigma\beta^{\mathcal C}\le
    M_{\mathcal C}\sigma^2_{\mathcal C}$ for a fixed $M_{\mathcal C}$, and $n\min_{j\in\Sz^{\mathcal C}}\Sigma_{jj}(\beta^{\mathcal
    C}_j)^2/(\sigma^2_{\mathcal C}\log n)\to\infty$.
\end{asmitems}
\end{assumption}

\begin{assumption}{NL}{No leakage, as seen by the start law}
Every null predictor $l\notin\Sz$ has a deterministic \emph{neighbourhood} $\Hl_l\supseteq\{l\}\cup\Sz$ with
$\abs{\Hl_l}\le e^{\kappa_H\sqrt{\log n}}$ for a fixed $\kappa_H$, such that, for every $U\not\ni l$ with $\abs{U}\le15$ and
$U_0:=U\cap\Hl_l$: (i) \emph{no leak:} $\abs{\Cov(W_l,W_j\mid W_{U_0})}\le C_Ln^{-1/2}$ for every $j\in\Sz\setminus U_0$,
with a fixed $C_L$; and (ii) \emph{locality:} $\abs{\Cov(W_l,W_v\mid W_{U_0})}\le r^\circ_n$ and
$\abs{\Cov(W_j,W_v\mid W_{U_0})}\le r^\circ_n$ for every $v\in U\setminus\Hl_l$ and $j\in\Sz\setminus U_0$, where
$r^\circ_n(n\log n)^{1/4}\to0$.
\end{assumption}

For a set $U$ and a predictor $l\notin U$: the \emph{signal gain} is $\delta^2_{U,l}:=\norm{(P_{U\cup\{l\}}-P_U)\mu}^2$;
the \emph{residual scale} is $\tau^2_U:=\sigs^2+\norm{(I-P_U)\mu}^2/(n-1)$; and $\hat\beta_{U,l}:=(X_U^\top X_U)^{-1}X_U^\top
x_l$, padded with zeros.

\begin{assumption}{reported-fixed}{Fixed-design visibility and leakage; A6 Assumption~A.13}
\begin{asmitems}
  \item\label{asm:rf-floor} \emph{Floors.} $\lmin(\Rhat_T)\ge\varphi_0^2$ for every $T$ with $\abs{T}\le16$ and for every
    $T=Q\cup\{l\}$, $l\in\set{p}$, with $Q$ as in~\ref{asm:rf-leak}.
  \item\label{asm:rf-vis} \emph{Visibility.} There is a deterministic $\omega_n\to\infty$ such that
    $\delta^2_{U,j}\ge\omega_n\log n$ for every $j\in\Sz$ and every $U\not\ni j$ with $\abs{U}\le15$.
  \item\label{asm:rf-leak} \emph{Bounded leakage.} There are a deterministic $Q\supseteq\Sz$ (with $\abs{Q}\le K$ and
    ($Q=\set{p}$ or $K\ge2\abs{Q}+32\sz+1$) for large $n$) and a fixed $\eta\in[0,\eta_A)$,
    $\eta_A:=(\sqrt{2A+1}-\sqrt{2A})^2$, such that $\delta^2_{U,l}\le\eta\,\tau^2_U\log n$ for every $l\notin Q$ and every
    $U\not\ni l$ with $\abs{U}\le15$.
  \item\label{asm:rf-size} \emph{Size of the leakage set.} One of: (d1) $\abs{Q\setminus\Sz}\le(1-\delta_Q)\log n$ for a
    fixed $\delta_Q\in(0,1)$; or (d2) $\abs{Q}\le n^b$ with $b<\tfrac12-\tfrac\eta2-\sqrt{2A\eta}$,
    $\theta^{*2}_{\min}n/(\abs{Q}+\log n)\to\infty$, and local regressions: for every $l$ a deterministic set
    $N_l\subseteq Q\setminus\{l\}$ with $\abs{N_l}\le h\log_2n$, for a fixed $h$ with $h+b<\tfrac12$, and a deterministic
    $L_n\to0$ with $L_n^2\log(\abs{Q}+1)=o(\log n)$, such that $\norm{\hat\beta_{T,l}-\hat\beta_{T\cap N_l,l}}_1\le L_n$
    for every $T\in\Lat(Q)$ with $\abs{T}<K$ and $l\notin T$.
\end{asmitems}
\end{assumption}

\begin{definition}[Explicit constants]\label{def:constants}
With $j_0:=\lceil(2A+1)\sz\rceil+1$ and $C_x:=Aj_0+\sz+\log(j_0+1)+2$:
$c_0(\sz,A):=\max\{4C_*,\,32[\sz(A+1)+2]\}$ with $C_*:=2\sz+3[Aj_0+\sz+\log(j_0+1)+1]=3C_x+2\sz-3$; and, with
$C_A:=8+3A$,
\[
  c_1(\sz,A):=\Big(\sqrt{2C_x}+\sqrt{2C_x+C_Aj_0+(2/\log2)\{\sz(A+1)+2\}}\Big)^2 .
\]
\end{definition}

\subsection{Conventions}\label{sec:conventions}

The conventions fix what ``the method'' means and idealise the computer; they are not statistical assumptions.

\begin{convention}{C1}{The method}
Exact arithmetic, the default configuration, valid inputs and seed. A run that meets it is \emph{admissible}; a sample
size $n$ is admissible if the method accepts some data set of that size as valid input. Statements hold for every valid
seed and every outcome of the calibration.
\end{convention}

\begin{convention}{C2}{Delivery}
\begin{enumerate}[label=(\roman*),leftmargin=2em]
  \item \emph{Idealised machine:} exact arithmetic at unit cost (including $\exp$ and $\log$); no integer, index,
    random-stream identifier or allocation overflow; wall time at most $c_{\mathrm t}\cdot(\text{operation count})+c_{\mathrm
    t0}$ and peak memory at most $c_{\mathrm m}\cdot(\text{words held})+c_{\mathrm m0}$ for fixed constants $(c_{\mathrm
    t},c_{\mathrm t0},c_{\mathrm m},c_{\mathrm m0})$; every allocation within the declared memory succeeds; the build and its
    dependencies work; no external interruption or worker crash.
  \item \emph{Covered configuration:} the default core of the software, with the local search from the start mode (cap
    $\ell_n+\abs{\tSz}$) and the rule that resolves ties in the start mode, implemented with
    $O\big(p\sum_k(2^{\abs{C_k}}+(K+1)^2)\big)$ operations and $O\big(\sum_k2^{\abs{C_k}}+J(K+1)\big)$ storage; the local
    search from the start mode starting at the smallest, lexicographically first maximiser of the start law; the default
    list of summaries only; certificate computations with declared polynomial operation and memory budgets, included in
    the allocation; serial execution or a fixed number of concurrent workers (memory envelope multiplied by that number);
    a fixed number of islands; the fallback temperature sequence switched on (the default).
  \item \emph{Random numbers and stops:} ideal independent uniform random numbers; the method's stop-and-return rule and
    stop-only resource guard as coded (a stopped run returns no estimates).
  \item \emph{Inputs and allocation:} a valid seed (an integer in $[1,2^{31}-1]$) and a declared memory and time
    allocation at least the method's explicit polynomial work and memory envelopes in $(n,\praw)$, including the memory
    checks made before the run and the bounds on the records and time of the local searches.
\end{enumerate}
\end{convention}

\begin{convention}{C2p}{Delivery, with the integer limits lifted}
An idealised machine that never overflows or fails (\as{C2}(i)), ideal random numbers, and the two fixed integer limits of Table~\ref{tab:params} (the ceiling
$10^6$ on $\ell_n$ and the bound $2^{31}-1$ on $\ell_n+\abs{\tSz}$) removed, everything else (other validators, the cap on the
enumeration used to find the start mode, and its fallback, every method cap) kept as coded, and an allocation of at least
the memory and time that the method's own polynomial requirement formulas require.
\end{convention}

\begin{convention}{C3}{Sampler}
\begin{enumerate}[label=(\roman*),leftmargin=2em,itemsep=1pt]
  \item The computer does exact arithmetic.
  \item Within each random stream, the random numbers are fresh, independent, ideal uniforms; the final run's stream for
    each island is independent of the data and of the calibration's streams.
  \item The coded start sampler returns exact draws from $\pi_0$, and the reported start normaliser equals the exact
    normalising constant of the unnormalised weights of $\pi_0$.
  \item The coded Markov moves realise their stated kernels, applied in a uniformly random order that is independent of
    the state.
  \item The outcome of the calibration is a function of the data and of the calibration's random streams, and each island's
    record is a measurable function of that outcome and of the island's stream.
  \item Statements about estimates refer to the potential (uninterrupted) run, which is identified with the reported
    values on completed runs.
\end{enumerate}
\end{convention}

\subsection{Which statements use which assumption}\label{sec:usage}

\begin{table}[ht]
\centering\small
\setlength{\tabcolsep}{4pt}
\caption{Assumptions and conventions against the statements that use them, read from the hypotheses of each statement.
Theorems~1--6 are in Section~\ref{sec:proved}; S-numbers are in Section~\ref{sec:supp}.}\label{tab:usage}
\begin{tabular}{@{}>{\raggedright\arraybackslash}p{0.3\linewidth}>{\raggedright\arraybackslash}p{0.62\linewidth}@{}}
\toprule
Assumption & Used by \\
\midrule
\USAGEROWS
\bottomrule
\end{tabular}
\end{table}
